% GENERATED FILE. Edit canonical lessons, then run npm run generate:books.
% canonical-source-hash: fnv1a64:0108da6427a6cc7c
% canonical-lessons: UR-C64-patta, UR-C64-pahar, UR-C64-jangal, UR-C64-darya, UR-C64-jhil

\chapter{Leaf, Mountain, Forest, River, Lake}
\label{ch:ur-leaf-mountain-forest-river-lake}
\hlchaptermodality{64}

\begin{chapteropening}
\textbf{By the end of this chapter:} \emph{I can say a leaf, a mountain, a forest, a river and a lake.}

Complete the last lesson of chapter 64: I can say a leaf, a mountain, a forest, a river and a lake.
\end{chapteropening}

\section[\texorpdfstring{\ur{پتا} (pattā)}{پتا (pattā)}]{\ur{پتا} (pattā) — a leaf}

\label{lesson:UR-C64-patta}



\begin{quote}
Before the new one: say the Urdu for a tree, then the Urdu for a flower.
\end{quote}

\subsection*{You'll want to know: \ur{پتا}}
\textbf{\ur{پتا}} — \emph{pattā} — \textquotedblleft{}a leaf\textquotedblright{}.

\begin{cousinweb}
\textbf{\ur{پتا}} (\emph{pattā}) is a tadbhava, a word inherited from Sanskrit and reshaped by sound change. It comes from \emph{pattra}, \textquotedblleft{}wing-feather; leaf\textquotedblright{} (CDIAL 7733).
\end{cousinweb}

\begin{grammarlens}[title={What you've built}]
One more word for this chapter.
\end{grammarlens}

\subsection*{Your turn}
\begin{itemize}
\raggedright
  \item \emph{Say it:} \emph{pattā}
  \item \emph{Say it:} \emph{pattā}, once more
  \item \emph{Say it:} say \emph{pattā}
  \item \emph{Recall:} say the Urdu for a tree, then the Urdu for a flower, then say \emph{pattā} again
\end{itemize}

\begin{tcolorbox}[breakable,colback=teal!4,colframe=teal!35!black,title={Before you move on}]
What does \ur{پتا} mean? (\textquotedblleft{}A leaf\textquotedblright{}.) Say it once more.
\end{tcolorbox}
\section[\texorpdfstring{\ur{پہاڑ} (pahāṛ)}{پہاڑ (pahāṛ)}]{\ur{پہاڑ} (pahāṛ) — a mountain}

\label{lesson:UR-C64-pahar}



\begin{quote}
Before the new one: say the Urdu for a flower, then the Urdu for a leaf.
\end{quote}

\subsection*{You'll want to know: \ur{پہاڑ}}
\textbf{\ur{پہاڑ}} — \emph{pahāṛ} — \textquotedblleft{}a mountain\textquotedblright{}.

\begin{cousinweb}
\textbf{\ur{پہاڑ}} (\emph{pahāṛ}) is a tadbhava, a word inherited and reshaped by sound change. It comes from \emph{*pāhāḍa}, \textquotedblleft{}rock, hill\textquotedblright{}, a form CDIAL reconstructs from its descendants (CDIAL 8141). CDIAL connects it, probably, with Sanskrit \emph{pāṣāṇa}, \textquotedblleft{}stone\textquotedblright{}.
\end{cousinweb}

\begin{grammarlens}[title={What you've built}]
One more word for this chapter.
\end{grammarlens}

\subsection*{Your turn}
\begin{itemize}
\raggedright
  \item \emph{Say it:} \emph{pahāṛ}
  \item \emph{Say it:} \emph{pahāṛ}, once more
  \item \emph{Say it:} say \emph{pahāṛ}
  \item \emph{Recall:} say the Urdu for a flower, then the Urdu for a leaf, then say \emph{pahāṛ} again
\end{itemize}

\begin{tcolorbox}[breakable,colback=teal!4,colframe=teal!35!black,title={Before you move on}]
What does \ur{پہاڑ} mean? (\textquotedblleft{}A mountain\textquotedblright{}.) Say it once more.
\end{tcolorbox}
\section[\texorpdfstring{\ur{جنگل} (jangal)}{جنگل (jangal)}]{\ur{جنگل} (jangal) — a forest}

\label{lesson:UR-C64-jangal}



\begin{quote}
Before the new one: say the Urdu for a leaf, then the Urdu for a mountain.
\end{quote}

\subsection*{You'll want to know: \ur{جنگل}}
\textbf{\ur{جنگل}} — \emph{jangal} — \textquotedblleft{}a forest\textquotedblright{}.

\begin{cousinweb}
\textbf{\ur{جنگل}} (\emph{jangal}) is a tadbhava, a word inherited from Sanskrit and reshaped by sound change. It comes from \emph{jāṅgala}, \textquotedblleft{}arid, sparse in trees\textquotedblright{} (CDIAL 5177). CDIAL notes a likely non-Aryan origin. The sense moved from \textquotedblleft{}waste land\textquotedblright{} to \textquotedblleft{}forest\textquotedblright{}, and English took it back as \emph{jungle}.
\end{cousinweb}

\begin{grammarlens}[title={What you've built}]
One more word for this chapter.
\end{grammarlens}

\subsection*{Your turn}
\begin{itemize}
\raggedright
  \item \emph{Say it:} \emph{jangal}
  \item \emph{Say it:} \emph{jangal}, once more
  \item \emph{Say it:} say \emph{jangal}
  \item \emph{Recall:} say the Urdu for a leaf, then the Urdu for a mountain, then say \emph{jangal} again
\end{itemize}

\begin{tcolorbox}[breakable,colback=teal!4,colframe=teal!35!black,title={Before you move on}]
What does \ur{جنگل} mean? (\textquotedblleft{}A forest\textquotedblright{}.) Say it once more.
\end{tcolorbox}
\section[\texorpdfstring{\ur{دریا} (daryā)}{دریا (daryā)}]{\ur{دریا} (daryā) — a river}

\label{lesson:UR-C64-darya}



\begin{quote}
Before the new one: say the Urdu for a mountain, then the Urdu for a forest.
\end{quote}

\subsection*{You'll want to know: \ur{دریا}}
\textbf{\ur{دریا}} — \emph{daryā} — \textquotedblleft{}a river\textquotedblright{}.

\begin{cousinweb}
\textbf{\ur{دریا}} (\emph{daryā}) is a loan, usually traced to Persian \emph{daryā}, \textquotedblleft{}sea, great river\textquotedblright{}.
\end{cousinweb}

\begin{grammarlens}[title={What you've built}]
One more word for this chapter.
\end{grammarlens}

\subsection*{Your turn}
\begin{itemize}
\raggedright
  \item \emph{Say it:} \emph{daryā}
  \item \emph{Say it:} \emph{daryā}, once more
  \item \emph{Say it:} say \emph{daryā}
  \item \emph{Recall:} say the Urdu for a mountain, then the Urdu for a forest, then say \emph{daryā} again
\end{itemize}

\begin{tcolorbox}[breakable,colback=teal!4,colframe=teal!35!black,title={Before you move on}]
What does \ur{دریا} mean? (\textquotedblleft{}A river\textquotedblright{}.) Say it once more.
\end{tcolorbox}
\section[\texorpdfstring{\ur{جھیل} (jhīl)}{جھیل (jhīl)}]{\ur{جھیل} (jhīl) — a lake}

\label{lesson:UR-C64-jhil}



\begin{quote}
Before the new one: say the Urdu for a forest, then the Urdu for a river.
\end{quote}

\subsection*{You'll want to know: \ur{جھیل}}
\textbf{\ur{جھیل}} — \emph{jhīl} — \textquotedblleft{}a lake\textquotedblright{}.

\begin{cousinweb}
\textbf{\ur{جھیل}} (\emph{jhīl}) is a tadbhava, a word inherited and reshaped by sound change. It comes from \emph{*jhilla}, \textquotedblleft{}swamp, lake\textquotedblright{}, a form CDIAL reconstructs from its descendants (CDIAL 5392).
\end{cousinweb}

\begin{grammarlens}[title={What you've built}]
One more word for this chapter.
\end{grammarlens}

\subsection*{Your turn}
\begin{itemize}
\raggedright
  \item \emph{Say it:} \emph{jhīl}
  \item \emph{Say it:} \emph{jhīl}, once more
  \item \emph{Say it:} say \emph{jhīl}
  \item \emph{Recall:} say the Urdu for a forest, then the Urdu for a river, then say \emph{jhīl} again
\end{itemize}

\begin{tcolorbox}[breakable,colback=teal!4,colframe=teal!35!black,title={Before you move on}]
What does \ur{جھیل} mean? (\textquotedblleft{}A lake\textquotedblright{}.) Say it once more.
\end{tcolorbox}
